\chapter*{\texorpdfstring{\textarabic{ملخص}}{Arabic Abstract}}
\addcontentsline{toc}{chapter}
  {\texorpdfstring{\textarabic{ملخص}}{Arabic Abstract}}

\begin{otherlanguage}{arabic}
تدعم قوائم مكونات البرمجيات شفافية سلسلة توريد البرمجيات وإدارة الثغرات
الأمنية، غير أن معظم مسارات إنشائها تعتمد على الشيفرة المصدرية أو ملفات
البناء أو بيانات مدير الحزم. وفي عمليات تدقيق تطبيقات أندرويد، قد لا يتوفر
سوى ملف \textenglish{APK} مترجم، حيث تندمج المكونات الخارجية داخل شيفرة
\textenglish{DEX} وتتعرض لتقليص الشيفرة وتحسينها وتمويهها.

تقدم هذه المذكرة الهندسية خط معالجة ساكنا، قائما على الأدلة ومحكوما بضوابط،
لاسترجاع قائمة مكونات البرمجيات من ملف \textenglish{APK}. يستخرج النظام
الأدلة من الأرشيف وملف التصريح والشيفرة والموارد والملفات الأصلية، ثم يقارن
خصائص التطبيق بملفات مرجعية لإصدارات مكتبات \textenglish{Maven}. كما يفصل
بين تحديد هوية المكون والثقة في إصداره الدقيق، ويحافظ على البدائل الملتبسة،
ويولد ملف \textenglish{CycloneDX JSON} مزودا بمعرفات
\textenglish{Package URL}، ثم يستعلم من خدمة \textenglish{OSV} مع مراعاة
درجة الثقة. وتستخدم أشجار تبعيات الشيفرة المصدرية لأغراض التقييم عند
توفرها، بينما تظل مساعدة التعلم الآلي و\textenglish{R8} محصورة في التقارير.

يفصل التقييم بين التطبيقات المضبوطة، والتطبيقات الحقيقية المتحقق منها
بالمصدر والمصنفة إلى نظيفة وجزئية، وتغطية المقام الكامل لتبعيات المصدر،
والتعرف ذي الحد الأعلى المقيد بإصدارات المصدر، والتجارب الفاشلة، وحالات
الحدود. استرجع النظام 84/84 إصدارا متوقعا في التطبيقات المضبوطة. ومن بين
102 تطبيق حقيقي متحقق منه بالمصدر أصدر مكونات \textenglish{Maven}، تأكدت 927/1,837 نسخة
\textenglish{Maven} صادرة، أي بنسبة تأكيد 50.5\% للإصدارات الصادرة، وليس
استدعاء لجميع التبعيات. كما توجد أربعة تطبيقات إضافية متحقق منها بالمصدر لم
تصدر أي مكون \textenglish{Maven}، وتبقى حالات ضابطة للاستدعاء. وبالمقارنة مع
12,082 سطرا لتبعيات محلولة من المصدر،
أصدر النظام 927 سطرا مطابقا من البداية إلى النهاية. كما أكد مسار مستقل مقيد
بإصدارات المصدر 6,772/6,772 سطرا صادرا، ولا يقدم هذا الرقم إلا بوصفه حدا
أعلى.

يثبت النموذج الأولي قابلية الاسترجاع العملي لقائمة مكونات البرمجيات انطلاقا
من الملف الثنائي، مع إظهار حدود تغطية قاعدة المراجع، واستدعاء المطابقة،
والتمييز بين الإصدارات المتجاورة وعائلات الملفات، والتعامل مع الشيفرة
الأصلية والمخفية، وقابلية التوسع، وحدود سلطة الذكاء الاصطناعي. ويتمثل مبدؤه
الأساسي في تفضيل معرفة جزئية قابلة للتدقيق مع امتناع صريح على ادعاء يقين غير
مدعوم بشأن الإصدار الدقيق.

\vspace{0.5cm}
\noindent\textbf{الكلمات المفتاحية ---}
أمن أندرويد، تحليل ملفات \textenglish{APK}، قائمة مكونات البرمجيات، كشف
المكتبات الخارجية، إسناد المكونات، \textenglish{CycloneDX}، إدارة الثغرات،
\textenglish{R8}.
\end{otherlanguage}
